\section{Harmonic lifting and two overlap estimates}\label{tra:harmonic-overlaps}

First-interaction deletion is effective at very small levels. Harmonic lifting gives a different overlap estimate whose loss is exponential only in the diagram deficit. It can then be interpolated with a dense entropy estimate.

This proof obtains dimension decay from two different overlap estimates. A cell-by-cell entropy argument controls the full regular trace with a subexponential error. A harmonic tuple decomposition supplies a bound with an error exponential only in the deficit. Interpolating the two local singular decompositions preserves a positive exponent through the recursion.

All logarithms are natural unless specified otherwise. We use the rotating-coordinate sweep of Section~\ref{sec:prelim}; it is denoted by $T_n$ here.

Use the following grid for the sweep. Split the bits into its first \(\lfloor d/2\rfloor\) and its remaining bits, and let \(m=2^{\lfloor d/2\rfloor}, q=n/m\). The \(q\) rows each have \(m\) positions (columns). Put
\[
          H=S_m^q,\qquad J=S_q^m
\]
for the row and column subgroups acting within rows and within columns, respectively. The first part of the sweep has independent \(T_m\)'s as components in \(H\), and the second has independent \(T_q\)'s in \(J\). Thus \(T_n\) is a product with those two factors (we can take \(T_1\) trivial).

\label{tra:harmonic-overlaps:setup}
Put $G=S_n$. All representation statements will be over \(\mathbb C\), with unitary models. By singular values in \(\lambda\) we mean in the irreducible \(V_\lambda\) of shape \(\lambda\); denote its dimension by \(D_\lambda\). We use Schatten norms (operator norm for \(p=\infty\)); matrix traces and Schatten norms use ordinary, not dimension-divided, trace. We use the following uniform exponent.

\begin{theorem}\label{tra:harmonic-overlaps:singular}
There is an absolute $a>0$ such that

\label{tra:harmonic-overlaps:all-index-conclusion}
\begin{equation}
 s_r(T_n\text{ in }\lambda)\ \le\ (D_\lambda r)^{-a},
 \qquad 1\le r\le D_\lambda,\quad n\text{ a power of }2,             \label{tra:harmonic-overlaps:eq1}
\end{equation}
with singular values ordered largest first. The claim as written allows 1 on a one-dimensional representation; on the nontrivial sign we can in fact use 0.

\end{theorem}

\label{tra:harmonic-overlaps:representation-and-hook-budgets}
We recall some standard results of representation theory that we use. We use symmetric group irreducibles indexed by partitions, with dimensions given by standard tableaux or the hook-length formula; transpose shapes have equal dimension. We use the content formula \cite[Eq.~(2.1), Proposition~5.3, Theorem~5.8]{VershikOkounkov2005} that the sum of all transpositions acts in shape \(\lambda\) as the sum of contents (column minus row) of its boxes. We also use the Pieri rule: in restriction to a product of two symmetric groups, the components trivial on the factor of size \(b\) are given on the other factor by removing a horizontal strip of size \(b\), with multiplicity one per such shape. Equivalently one can use induction with the trivial representation of the indicated factor. In particular, invariance under the subgroup of independent switches at the first layer requires a shape obtainable by adding \(n/2\) horizontal strips of size 2, so
\begin{equation}
                 T_n\text{ in }\lambda=0\quad\text{if }\lambda'_1>n/2,      \label{tra:harmonic-overlaps:eq2}
\end{equation}
where prime denotes transpose.

For \(\mu\vdash M\), set \(j=M-\mu_1\), the deficit, and let \(\bar\mu\) be the shape below the top row. Write
\[
 B_\mu = D_\mu/D_{\bar\mu}.
\]
The hook formula, taking the contribution of just the top row, gives
\begin{equation}
                      2^{-j}\binom Mj\ \le B_\mu\le \binom Mj.        \label{tra:harmonic-overlaps:eq3}
\end{equation}
Indeed the factor by which the hook product along that row exceeds \(\mu_1!\) is at most \(2^j\), by the numbers of boxes below the boxes along the row. Also throughout we can use
\begin{equation}
                        (M)_t/M^t\ \ge\ e^{-C t}\qquad (0\le t\le M),      \label{tra:harmonic-overlaps:eq4}
\end{equation}
with \((M)_t=M(M-1)\cdots(M-t+1)\); for example average the negative logs of the factors, bounding by the average for \(t=M\).

We will work with projections specified locally in the grid as follows. In one row take an orthogonal projection in the matrix algebra of one specified irreducible \(\mu\vdash m\), of rank \(l\), extended by zero on the other irreducibles. We may suppose projections are nonzero. Regard it as a group algebra projection on the row group. Its rank in the regular representation of that group is \(D_\mu l\). Suppose we choose such projections in every row, which tensor together to give \(P_H\), and similarly choose a projection \(P_J\) using columns. Use \(R_H\), \(R_J\) to denote products of the regular ranks of the local projections on the respective sides, or upper bounds given by products of local upper bounds. We study the overlap \(P_H P_J\), including in individual irreducibles of \(G\).

\subsection{A regular overlap from cell entropy}

\begin{proposition}\label{tra:harmonic-overlaps:dense}

\label{tra:harmonic-overlaps:dense-overlap-statement}
For the balanced sizes \(m,q\) here, let \(L=\log n\), and for sufficiently large \(n\) write \(\epsilon=L^{-1/5}\). We claim
\begin{equation}
       \|P_H P_J\|_{4,\mathrm{reg}(G)}^4
              \le (R_H R_J)^{1+\epsilon}
                       \exp\{ n\exp(-L^{1/2})\}.                            \label{tra:harmonic-overlaps:eq5}
\end{equation}
This estimate in particular bounds the Schatten norm within shape \(\lambda\), with the fourth power on the right divided by \(D_\lambda\), by the multiplicity in the regular representation.

\end{proposition}\begin{proof}\label{tra:harmonic-overlaps:dense-proof}

\label{tra:harmonic-overlaps:coefficients-and-tilted-entropy}
Write the ordinary group algebra coefficients of the projections as
$p(h),w(y)$ on $H,J$, so their sums give the operators without a
uniform-measure prefactor. Lemma~\ref{lem:coefficient-compatibility},
applied to the local projections in this normalization, bounds
$\|P_HP_J\|_{4,\mathrm{reg}(G)}^4$ by
\begin{equation}
                   n! \sum_{(h,y):\, (hy)^{-1}\in HJ}|p(h)|^2 |w(y)|^2.   \label{tra:harmonic-overlaps:eq6}
\end{equation}
The condition has a useful description. In the grid of middle positions in a product \(hy\), the column permutation step leading up to a middle position gives a source-row symbol, and the row permutation step from it gives a target-column symbol. Thus there is one column symbol at each cell \((i,b)\), with these symbols a permutation of \([m]\) along each row \(i\); and there is one row symbol there, a permutation of \([q]\) along each column \(b\). The condition is that the pairs of symbols cover the \(m q=n\) possibilities once each. In fact, to go back by another product in \(HJ\), each target column must send the ending row indices to the prescribed source rows by a permutation, which requires and, for this part, is ensured by this pair condition. Then within any resulting source row we get the prescribed permutation to restore the source columns as well.

The same lemma normalizes the squared line coefficients as independent
probability laws whose densities relative to uniform are bounded by
the local regular ranks. Consequently \eqref{tra:harmonic-overlaps:eq6}
is at most
\begin{equation}
       \frac{n! R_H R_J}{|H||J|}\,\Pr(\text{pair covering condition}),     \label{tra:harmonic-overlaps:eq7}
\end{equation}
where permutations across rows and columns are independent with tilted laws, each having density against uniform bounded by the local regular rank (we can use its upper bound; the laws come from the normalized squares). Inverses of permutations in setting up the array of symbols are immaterial.

Here is the needed estimate for the probability. Suppose it is positive, and consider the law of the entire array conditioned on pair covering; in the entropy argument this law already includes the tilts. Let \(\mathcal E\) denote entropy, using \(X\) for the array; let \(X_i\) denote the vector of column symbols along row \(i\), and \(Y_b\) the vector of row symbols down column \(b\). For a density against uniform, its expected log at any law is bounded above by that law's relative entropy against uniform. It is also bounded above by the log of the sup of the density. Using weights \(1-\epsilon,\epsilon\) on these two bounds, writing the log of the probability from \eqref{tra:harmonic-overlaps:eq7} by entropy of the conditioned law gives
\begin{equation}
\begin{split}
 \log\Pr(\text{pair covering})
 \ \le\ &\mathcal E(X)
   -(1-\epsilon)\Big(\sum_i\mathcal E(X_i)+\sum_b\mathcal E(Y_b)\Big)\\
 &-\epsilon\log(|H||J|) + \epsilon\log(R_H R_J).                        \end{split}
\label{tra:harmonic-overlaps:eq8}
\end{equation}

\label{tra:harmonic-overlaps:random-cell-priorities}
We bound the entropy difference here by exposing cells in random order. The following details, in particular keeping the marginal entropies, help make the rank power loss in \eqref{tra:harmonic-overlaps:eq5} small. Give every cell an independent uniform \([0,1]\) priority, exposing in decreasing order. At cell \(e=(i,b)\), for given order, use the conditional distributions
\begin{itemize}
\item \(u\) of its column symbol given just the previous column symbols along row \(i\);
\item \(v\) of its row symbol given just the previous row symbols along column \(b\).

\end{itemize}
These are for the marginal laws appearing in \eqref{tra:harmonic-overlaps:eq8}. They are supported on at most \(M_e,Q_e\) choices respectively, the counts not already used in that row or column, including the current symbol. By the entropy chain rule and the variational bound for log sums, the contribution in the entropy difference from exposing this cell is at most the expectation of
\begin{equation}
            \log\sum_{(c,s)\ \mathrm{available}} u(c)^{1-\epsilon}
                                                       v(s)^{1-\epsilon},   \label{tra:harmonic-overlaps:eq9}
\end{equation}
where we also require that the pair has not been used in any previous cell. Indeed the conditional array law uses only available possibilities, and the negative marginal entropy contributions are computed exactly by averaging logs of \(u,v\).

The unrestricted sum divided by \((M_e Q_e)^\epsilon\) is at most 1. Condition on the full array assignment, the priority \(x\) of cell \(e\), and priorities in its row or column. The distributions \(u,v\) depend, given the assignment, only on the respective in-row, in-column previous data; they do not depend on the remaining priorities. Take the cell owning each pair in this full assignment. Owners outside row \(i\) and column \(b\) leave their pair still unused with probability \(x\). The contribution of owners inside row \(i\) to the normalized sum is at most \(Q_e^{-\epsilon}\), since each \(c\) has only one such pair there. That of owners in column \(b\) is at most \(M_e^{-\epsilon}\). Therefore Jensen bounds the conditional expectation of \eqref{tra:harmonic-overlaps:eq9} minus \(\epsilon\log(M_e Q_e)\) by
\[
                        \log(x+Q_e^{-\epsilon}+M_e^{-\epsilon}).
\]
We can always replace a positive bound on this difference by 0.

Averaging this last estimate gives, uniformly as \(n\) gets large, at most
\[
                              -1+O(\exp(-L^{3/5})).
\]
To see this, when \(x\ge x_0=\exp(-L^{2/3})\), except with probability smaller than \(\exp(-n^{c_0})\) for some \(c_0>0\), both remaining counts are at least \(n^{1/5}\), by usual binomial bounds. Thus on the range \(x\ge x_0\) the result follows by comparison with the integral of \(\log x\), since \(2\exp(-L^{4/5}/5)\) is sufficiently small; the negligible binomial error can use the bound 0, as can \(x<x_0\), with the missing integral there of order \(x_0(1+|\log x_0|)\).

Summing \eqref{tra:harmonic-overlaps:eq9}, the terms \(\epsilon\log(M_e Q_e)\) sum to \(\epsilon\log(|H||J|)\) in any order. We obtain from \eqref{tra:harmonic-overlaps:eq8}
\[
 \log\Pr(\text{pair covering})\le
           -n+O(n\exp(-L^{3/5}))+\epsilon\log(R_HR_J).
\]
Finally \(n!/(|H||J|)\le \exp(n+O(\log n))\) by factorial estimates. This proves \eqref{tra:harmonic-overlaps:eq5}.

\end{proof}

\subsection{Lifting harmonic functions on injections}

For the second overlap estimate we use representations on ordered distinct tuples. Here and elsewhere functions on finite sets are given uniform-measure norms unless indicated.

\label{tra:harmonic-overlaps:injection-definition}
For $0\le t\le M$, let \(\mathcal I_{M,t}\) be functions on injections of \(t\) ordered slots into \([M]\), with label action by \(S_M\) (in application the labels \([M]\) here are positions). Functions of only a subset of the slots lift isometrically to \(\mathcal I_{M,t}\). Say a function is harmonic here if it sums to zero when varying any one slot with the others fixed. Equivalently the projection on ignoring any one slot is zero. Its extension by zero on repetitions is centered in every slot even under independent uniform sampling with replacement.

\begin{lemma}\label{tra:harmonic-overlaps:lifting}

\label{tra:harmonic-overlaps:lifting-statement}
The harmonic part of \(\mathcal I_{M,t}\) consists exactly of the shapes \(\mu\vdash M\) with \(M-\mu_1=t\), each with multiplicity \(D_{\bar\mu}\). More generally an occurring shape has \(j=M-\mu_1\le t\). Also, in the component of such a shape, if we write functions on \(t\) slots as sums of lifts from just \(j\) slots, we can take coefficient functions from the same shape with sum of squared norms bounded by \(e^{C t}\) times the squared norm on \(t\) slots, using a linear choice of coefficients. The coefficient selection respects group algebra projections acting on the \(S_M\) shape.

\end{lemma}\begin{proof}\label{tra:harmonic-overlaps:lifting-proof}

\label{tra:harmonic-overlaps:deletion-spectrum}
We first identify the harmonic part through the spectrum of coordinate deletion, then bound the lifting loss uniformly over the full range \(0\le t\le M\).

The tuple representation is that on cosets of \(S_{M-t}\), with commuting slot action. Multiplicity space for \(\mu\) is given, up to dual convention, by the \(S_{M-t}\)-invariants, and under the slot group \(S_t\) it has types \(\nu\vdash t\) with \(\mu/\nu\) a horizontal strip of size \(s=M-t\), by Pieri, each once. This in particular requires \(j\le t\). Let \(Q_t\) be the sum of projections ignoring one slot, over the \(t\) slots. It acts within \(\mu,\nu\) with eigenvalue
\begin{equation}
                 \frac{t+\mathrm{ct}(\mu)-\mathrm{ct}(\nu)-\binom s2}
                      {s+1},                                             \label{tra:harmonic-overlaps:eq10}
\end{equation}
where \(\mathrm{ct}\) is the content sum. In fact in the coset model each projection averages the identity and swapping the slot with any of the \(s\) unused ones; the cross-transposition sum in \(Q_t\) is the full transposition sum minus those of the two sets of slots.

For strip columns \(c\), using one-based column numbers, put
\[
                             p=\sum_{\rm strip} c-s(s+1)/2\ge 0.
\]
Then \eqref{tra:harmonic-overlaps:eq10} is at least
\begin{equation}
                             (t-j+p)/(s+1),                              \label{tra:harmonic-overlaps:eq11}
\end{equation}
since the contribution to the content difference in each such column is \(c-\mu'_c\), and the total excess of heights over 1 is \(j\). Thus \(Q_t>0\) on the indicated components with \(t>j\). When \(j=t\), the shape does not occur if we drop a slot, and here the strip removes one box in every column so the multiplicity is \(D_{\bar\mu}\). This proves the statement about harmonic functions (also immediate at \(t=0\)).

\label{tra:harmonic-overlaps:lifting-product-bound}
Within shape \(\mu\) in \(\mathcal I_{M,u}\), let \(L_S\) be the isometric lift from the slots in \(S\), and put
\[
 \mathcal A_u=\sum_{|S|=j}L_SL_S^*,\qquad
 a_u=\lambda_{\min}(\mathcal A_u).
\]
Here \(j\le u\le t\), and \(\mathcal A_j=I\). For \(u>j\), let \(L_i\) lift from all slots except \(i\), and let \(\mathcal A_{u-1}^{(i)}\) be the corresponding operator on those slots. Each \(j\)-subset omits exactly \(u-j\) slots, so
\[
 (u-j)\mathcal A_u
   =\sum_{i=1}^u L_i\mathcal A_{u-1}^{(i)}L_i^*
   \ \ge\ a_{u-1}\sum_{i=1}^u L_iL_i^*
   =a_{u-1}Q_u.
\]
All lifts preserve the label type \(\mu\). Thus \(a_u\) is bounded below by \(a_{u-1}\) times the minimum of \eqref{tra:harmonic-overlaps:eq10} at slot count \(u\), divided by \(u-j\). We show that the product of these lower factors, capped at 1 if desired, is at least \(\exp(-Ct)\).

Here is the product check. Shape \(\mu\) stays fixed. Write \(w=M-j\), the width, and at a slot count \(u\) with \(j<u\le t\), write \(s_u=M-u\), \(l=u-j=w-s_u\). By \eqref{tra:harmonic-overlaps:eq11}, the lower factor is at least
\begin{equation}
                    (l+p)/((s_u+1)l),                                    \label{tra:harmonic-overlaps:eq12}
\end{equation}
taking the worst allowable strip, now of size \(s_u\).
If \(t<M/4\), the last band of columns of \(\mu\), of height 1, has length \(b\ge M-2j\). In any band of equal height removed columns must form a suffix. Thus the sum of the \(l\) unremoved column numbers is at most the sum from taking the first \(l\) columns of the last band (using \(l\le b\)). Consequently
\[
 p\ge l(b-l),\qquad
 \frac{l+p}{(s_u+1)l}
   \ge\frac{1+b-l}{s_u+1}
   \ge1-\frac{j}{s_u+1}>\frac23.
\]
The product therefore loses at most a constant to the power \(t\).

If \(t\ge M/4\), we bound the total log loss by \(O(M)=O(t)\). View each size \(s_u\) as a cut after column \(s_u\), and let \(b_r\) be the band lengths, indexed by height \(r\). If a cut lies inside a band, let \(v\) be its distance to the nearer boundary; at a boundary put \(v=0\). Compare a permissible strip with the first \(s_u\) columns. The number \(h\) of missed columns before the cut equals the number of selected columns after it. Within the band the selected columns form a suffix, so either all columns before the cut are missed or all columns after it are selected. Hence \(h\ge v\). Even the cheapest exchange of \(h\) columns across the cut increases the column sum by \(h^2\), giving
\[
 p\ge h^2\ge v^2,\qquad
 \frac{l+p}{(s_u+1)l}\ge\frac{(1+v)^2}{4M^2}.
\]
A bound for the logarithmic loss in \eqref{tra:harmonic-overlaps:eq12} is therefore
\[
                         2\log(CM/(1+v)).
\]
Summing within each band, including its endpoints, gives \(O(\sum_r b_r\log(eM/b_r))\) by a factorial estimate. To bound this sum, put \(\pi_r=b_r/w\). This probability distribution on positive heights has mean \(M/w\). Comparison with the geometric law of that mean gives its entropy \(\mathcal E(\pi)=-\sum_r\pi_r\log\pi_r\le\log(eM/w)\). Therefore
\[
 \sum_r b_r\log\frac{eM}{b_r}
 =w\log\frac{eM}{w}+w\mathcal E(\pi)
 \le2w\log\frac{eM}{w}\le2M.
\]
Since \(M\le4t\), the product has the required lower bound \(\exp(-Ct)\).

Let \(\mathcal L\) add the lifts from the direct sum of the \(j\)-slot spaces to the \(t\)-slot space, within shape \(\mu\). Then \(\mathcal L\mathcal L^*=\mathcal A_t\), so \(\mathcal L^*\mathcal A_t^{-1}\) selects coefficients linearly, with squared norm at most \(a_t^{-1}\le\exp(Ct)\). This map is label-equivariant and hence preserves the image of every group algebra projection on the shape. This proves the lifting facts.

\end{proof}

\subsection{A sparse-deficit overlap}

\begin{proposition}\label{tra:harmonic-overlaps:sparse-overlap}

\label{tra:harmonic-overlaps:sparse-overlap-statement}
For every \(\lambda\vdash n\) with deficit \(k=n-\lambda_1\ge 1\),
\begin{equation}
              \|P_H P_J\|_{4,\lambda}^4
                    \le \exp(C k)\,\frac{R_H R_J}{D_\lambda}.             \label{tra:harmonic-overlaps:eq13}
\end{equation}
\end{proposition}

The loss here is exponential in \(k\). For large deficits we will instead use the regular-trace estimate \eqref{tra:harmonic-overlaps:eq5}.

\begin{proof}\label{tra:harmonic-overlaps:sparse-overlap-proof}

\label{tra:harmonic-overlaps:fine-slot-expansion}
Consider the harmonic representation in \(\mathcal I_{n,k}\). It contains \(D_{\bar\lambda}\) copies of \(\lambda\), so suffices for getting the estimate with this multiplicity taken into account. Both projections preserve this representation. We estimate their angle operator: that between their ranges, with entries given by inner product, whose singular values are those for the product apart from zeros.

Tuples will have row and column vectors \(a=(a_i)_{i=1}^k\), \(b=(b_i)_{i=1}^k\) respectively, and put \(D(a,b)=1\) if the cells are distinct and 0 otherwise. We compare norms and probabilities to the model with \(a,b\) independent uniform, including across slots (called iid sampling below). Within row \(r\) suppose the specified shape of \(P_H\) is \(\mu_r\), of deficit \(j_r\), with local projection rank \(l_r^*\) in that shape. Fixing \(a\), let \(t_r\) be the number of slots in row \(r\). On a tuple fiber with this assignment the row group acts just on the corresponding injection of \(t_r\) columns. For \(P_H\) to be nonzero here we need \(j_r\le t_r\); in particular we can assume
\begin{equation}
                                \sum_r j_r\le k                           \label{tra:harmonic-overlaps:eq14}
\end{equation}
(and analogously in columns).

For \(f\) in the harmonic representation in the range of \(P_H\), we can therefore expand on distinct tuples into lifts of the form
\begin{equation}
                                  F_S(a,b_S),                             \label{tra:harmonic-overlaps:eq15}
\end{equation}
indexed by sets \(S\) of slots of size \(\sum_r j_r\). Each term only uses assignments \(a\) whose restriction \(a_S\) has exactly \(j_r\) slots in each row; outside such assignments take it as zero. As a function of \(b_S\), it uses the projection range on the respective \(j_r\) fine slots over each row, in the corresponding shape there. In particular that range, taking all these rows together and extending by zero off distinct column values in the same row among the indicated slots, is a space \(W(a_S)\) of rank
\begin{equation}
                 T_H=\prod_r l_r^* D_{\bar\mu_r}
                         \le R_H \Big/\prod_r B_{\mu_r}.                  \label{tra:harmonic-overlaps:eq16}
\end{equation}
We will give this fine-slot space the iid \(b_S\) norm. It depends only on the assignments \(a_S\). It is the natural relabeling to those slots of a tensor product of row spaces, each invariant under slot permutations there (group algebra on positions commutes with slot permutations).

The expansion \eqref{tra:harmonic-overlaps:eq15} follows by the injection lifting fact, used given the row assignment, with the coefficient maps tensored within that assignment. Ranging over assignments and using iid norms, the terms are obtained linearly with
\begin{equation}
                        \sum_S \|F_S\|_{\rm iid}^2
                                     \le \exp(Ck)\|f\|^2.                 \label{tra:harmonic-overlaps:eq17}
\end{equation}
For this, conditional tuple measure on a possible assignment gives independent uniform injections over the row fibers, and assignment weights are comparable to uniform iid weights up to \(\exp(Ck)\) on possible assignments: the ratio is given by the falling factorials within rows and the overall one, so we can use \eqref{tra:harmonic-overlaps:eq4}. We use zero on impossible assignments; extension by zero for the fine-slot functions also respects the norm bound. The product of the lifting bounds over all rows costs only \(\exp(C\sum_r t_r)=\exp(Ck)\).

Similarly expand harmonic \(g\) in the range of \(P_J\) into lifts
\begin{equation}
                                  G_T(b,a_T),                              \label{tra:harmonic-overlaps:eq18}
\end{equation}
where, using letters \(c\) for columns, column shapes are \(\gamma_c\) of deficits \(d_c^*\) and \(T\) has \(\sum_c d_c^*\le k\) slots. Required \(b_T\) assignments use these counts. The analogous fine-slot space \(W'(b_T)\) has rank \(T_J\le R_J/\prod_c B_{\gamma_c}\).

\label{tra:harmonic-overlaps:centering-and-residual-coverage}
We now express the overlap through the coefficient spaces of these expansions. For each \(S\), let \(\mathcal H_S\) consist of functions \(F(a,b_S)\) whose fine section lies in \(W(a_S)\), with zero values unless \(a_S\) has the required row counts. Give it the norm
\[
 \|F\|_{\mathcal H_S}^2
   =\mathbb E_a\mathbb E_{b_S}|F(a,b_S)|^2
\]
under iid sampling. Define \(\mathcal K_T\) analogously for \(E(b,a_T)\) with fine sections in \(W'(b_T)\). Outside the marked slots, the coarse assignments are unrestricted. Orthonormal fine-slot bases identify these norms with Euclidean coefficient norms averaged over the coarse vectors.

Let \(A_Sf=F_S\) and \(B_Tg=G_T\) be the linear coefficient selections above, on the harmonic ranges of \(P_H\) and \(P_J\), respectively. Equation~\eqref{tra:harmonic-overlaps:eq17} and its column analogue give
\[
 \sum_S\|A_Sf\|_{\mathcal H_S}^2\le e^{Ck}\|f\|^2,
 \qquad
 \sum_T\|B_Tg\|_{\mathcal K_T}^2\le e^{Ck}\|g\|^2.
\]
Their individual operator norms are therefore at most \(e^{Ck}\), after adjusting the constant.

Write \(\mathsf C_U=\prod_{i\in U}(\mathrm{Id}-\mathbb E_{b_i})\). The zero extension of harmonic \(f\) is centered in every whole cell-slot under iid sampling. Against it, we may first center each \(G_T(b,a_T)\) in the slots outside \(T\). Since \(G_T\) is independent of \(a_i\) there, this applies exactly \(\mathsf C_{[k]\setminus T}\). Now expand \(f\) on distinct tuples. For each \(S,T\), put
\[
                   I=[k]\setminus (S\cup T),\qquad i_0=|I|.
\]
The independence of \(F_S\) from \(b_I\) and self-adjointness of centering give
\[
 \mathbb E_{\rm iid}
   [D\,\overline{F_S}\,\mathsf C_{[k]\setminus T}G_T]
 =\mathbb E_{\rm iid}
   [C_I\,\overline{F_S}\,\mathsf C_{S\setminus T}G_T],
 \qquad C_I=\mathsf C_I D.
\]
The remaining centering acts only on coarse coordinates outside \(T\). It is a contraction on \(\mathcal K_T\) and leaves \(W'(b_T)\) unchanged.

Define \(\mathcal M_{S,T}:\mathcal K_T\to\mathcal H_S\), abbreviated to \(\mathcal M\) when the pair is fixed, by
\begin{equation}
 \langle F,\mathcal M E\rangle=
                        \mathbb E_{\rm iid}
                  \big[C_I(a,b)\overline{F(a,b_S)}E(b,a_T)\big].          \label{tra:harmonic-overlaps:eq19}
\end{equation}
In orthonormal fine-slot bases, its kernel entries from coarse \(b\) to coarse \(a\) are \(C_I(a,b)\) times the conjugate and unconjugated fine basis evaluations at \(b_S,a_T\). With \(c_k=n^k/(n)_k\), the angle operator between the two harmonic projection ranges is therefore
\begin{equation}\label{tra:harmonic-overlaps:angle-factorization}
 \mathcal O=c_k\sum_{S,T}
       A_S^*\mathcal M_{S,T}\mathsf C_{S\setminus T}B_T.
\end{equation}
Here \(c_k\le e^{Ck}\), and there are at most \(4^k\) pairs \(S,T\).

The purpose of the two norm estimates below is the inequality
\[
 \|\mathcal M\|_4^4
 \le \|\mathcal M\|_{\rm op}^2\|\mathcal M\|_{\rm HS}^2.
\]
The residual slots will supply a factor $(k/n)^{i_0/2}$ in each squared norm. Counting the marked slots will supply the remaining factor $(k/n)^{|S\cup T|}$ in the Hilbert--Schmidt estimate, together with $R_HR_J$. Since $i_0+|S\cup T|=k$, their product gives $(k/n)^k$. Finally, the $D_{\bar\lambda}$ copies of $\lambda$ in the harmonic representation convert this estimate to the $D_\lambda$ denominator in \eqref{tra:harmonic-overlaps:eq13}. All coefficient maps and subset sums will cost only $\exp(Ck)$.

First, with \(\delta=k/n\), uniformly given the values in all marked slots \(S\cup T\), there is the residual estimate
\begin{equation}
           \mathbb E_{a_I,b_I}|C_I(a,b)|^2
                              \le \exp(Ck)\delta^{i_0/2}.                 \label{tra:harmonic-overlaps:eq20}
\end{equation}
For clarity, expand the centering operators using independent alternatives to \(b_i,\ i\in I\). In the alternating differences, if a slot \(i\in I\) cannot create a duplicate cell with some other slot (allowing the displayed and alternative values in any slots with alternatives), the difference sum is zero. The squared quantity can thus be bounded with a factor at most \(4^k\) by the probability of the potential-duplicate coverage event, using also Jensen over alternatives. The residual slots are sampled independently, with their alternatives (same row in the two values for the slot). In the collision graph every one of these slots must be incident to an edge. To bound the probability, in components with a connection to a marked slot use edges of a forest leading to marked roots, and in other components use a forest leading to a free root; there are at least two slots in any free component. If the number of edges chosen is \(e\), then \(e\ge i_0/2\). The choices are bounded by \(2^{O(k)}k^e\), for example by giving nonroots parents. Ordered from the roots, edge conditions each have probability at most \(4/n\) for a new slot, since each allowed cell-value version in it is uniform. A union bound proves \eqref{tra:harmonic-overlaps:eq20}, also when the desired estimate is loose because \(\delta\) is not small.

\label{tra:harmonic-overlaps:operator-bound-and-overlap-assignments}
Fix all marked cell values. In the residual variables, \(F\) depends only on \(a_I\) and \(E\) only on \(b_I\), so
\[
 \left|\mathbb E_{a_I,b_I}[C_I\overline FE]\right|
 \le\left(\mathbb E_{a_I,b_I}|C_I|^2\right)^{1/2}
     \left(\mathbb E_{a_I}|F|^2\right)^{1/2}
     \left(\mathbb E_{b_I}|E|^2\right)^{1/2}.
\]
Apply \eqref{tra:harmonic-overlaps:eq20}, then average over the marked values and use Cauchy--Schwarz once more. The resulting norms are exactly \(\|F\|_{\mathcal H_S}\) and \(\|E\|_{\mathcal K_T}\), since marked coordinates absent from a function integrate out. Hence
\begin{equation}
                         \|\mathcal M\|_{\rm op}^2
                                  \le \exp(Ck)\delta^{i_0/2}.             \label{tra:harmonic-overlaps:eq21}
\end{equation}
The full Hilbert-Schmidt bound takes more accounting for marks. In a normalized finite sampling measure it is the expected squared kernel size (sum of squared entries in the bases), thus uses
\begin{equation}
                 |C_I(a,b)|^2\, U(a_S,b_S)\, V(b_T,a_T),                 \label{tra:harmonic-overlaps:eq22}
\end{equation}
where \(U,V\) are the sums of squared basis evaluations of the fine-slot spaces, set to zero at assignments not having the required counts. After integrating out residual slots we are bounded using \eqref{tra:harmonic-overlaps:eq20}, leaving the expectation involving just \(U,V\).

Write
\[
                    O=S\cap T,\quad z=|O|,\quad
                    x=|S\setminus T|,\quad y=|T\setminus S|.
\]
We evaluate this expectation by fixing the \(a_S,b_T\) assignments and averaging the fine evaluations in the non-overlap coordinates \(b_{S\setminus T}, a_{T\setminus S}\). Suppose the counts of the overlap slots in those assignments are \(h_r\le j_r\) in rows and \(v_c\le d_c^*\) in columns. The numbers of non-overlap assignments, given overlap assignments, are
\begin{equation}
                       \frac{x!}{\prod_r(j_r-h_r)!},\qquad
                       \frac{y!}{\prod_c(d_c^*-v_c)!}.                   \label{tra:harmonic-overlaps:eq23}
\end{equation}
The averaged \(U\) value is \(T_H m^z\) times the probability of the evaluation values \(b_O\) in a probability distribution on those slots. Indeed \(U/T_H\) is a probability density on \(b_S\) against iid measure. Moreover this is a product over rows, supported on distinct values per row before and after marginalization, with the overlap distribution permutation invariant within each row. And per row it is always the same distribution up to slot naming for the given \(h_r\), regardless of which overlap slots were assigned or how the assignment was extended to the non-overlap slots. These properties follow from the tensor description and slot invariance of the spaces giving sums of squared basis evaluations. Similarly the averaged \(V\) is \(T_J q^z\) times a probability of \(a_O\), with the analogous properties over columns.

In summing over overlap assignments \(a_O,b_O\) at fixed counts \(h_r,v_c\), the sum of products of these two probabilities is at most
\begin{equation}
                             \frac{z!}{\prod_r h_r!\prod_c v_c!}.           \label{tra:harmonic-overlaps:eq24}
\end{equation}
To verify this, each contributing overlap with positive products forms a simple bipartite graph between the rows and columns with edges labeled by overlap slots: no repeated cell is allowed by the supports. A given simple graph with these degrees has \(z!\) labelings. The row-side probability of the ordered column values is the probability of the corresponding unordered neighbor choices independently by rows, divided by \(\prod h_r!\); likewise on the other side. Summing products of the unordered probabilities over such graphs costs at most 1.

\label{tra:harmonic-overlaps:hilbert-schmidt-and-fourth-norm}
The normalization for \(a_S,b_T\) is \(q^{-(x+z)}m^{-(y+z)}\), which with the \(m^z q^z\) above leaves \(q^{-x}m^{-y}\). Thus \eqref{tra:harmonic-overlaps:eq22}, after all averaging, is bounded by a sum over overlap counts \(h_r,v_c\) using terms at most
\begin{equation}
 \exp(Ck)\delta^{i_0/2} T_H T_J
       q^{-x}m^{-y}
       \frac{x!y! z!}
        {\prod_r (j_r-h_r)! h_r!\,\prod_c(d_c^*-v_c)! v_c!}.             \label{tra:harmonic-overlaps:eq25}
\end{equation}
There are at most \(\exp(Ck)\) count choices, using \eqref{tra:harmonic-overlaps:eq14}. Use \eqref{tra:harmonic-overlaps:eq16} and its analogue, with
\( B_{\mu_r}\ge \exp(-Cj_r)m^{j_r}/j_r!\)
by \eqref{tra:harmonic-overlaps:eq3}, \eqref{tra:harmonic-overlaps:eq4}, and likewise in columns. Since \(\sum j_r=x+z,\ \sum d_c^*=y+z\), \eqref{tra:harmonic-overlaps:eq25} shows
\begin{equation}
                 \|\mathcal M\|_{\rm HS}^2
                     \le \exp(Ck) R_H R_J\,\delta^{i_0/2+x+y+z}.         \label{tra:harmonic-overlaps:eq26}
\end{equation}
Here \(x! y! z!\le k^{x+y+z}\), and the split factorials within rows or columns versus their combined factorials only cost \(\exp(Ck)\).

Now \eqref{tra:harmonic-overlaps:eq21}, \eqref{tra:harmonic-overlaps:eq26} give
\[
                   \|\mathcal M\|_4^4\le
                         \exp(Ck) R_H R_J\,\delta^k.
\]
Apply the Schatten ideal and triangle inequalities to \eqref{tra:harmonic-overlaps:angle-factorization}. The coefficient-map bounds, centering contractions, injection normalization and at most \(4^k\) pairs give
\[
 \|\mathcal O\|_4^4
 \le e^{Ck}\max_{S,T}\|\mathcal M_{S,T}\|_4^4
 \le e^{Ck}R_HR_J\delta^k.
\]
This bounds the angle operator in the harmonic representation, which contains \(D_{\bar\lambda}\) copies of \(\lambda\). Since \(D_\lambda/D_{\bar\lambda}\le(en/k)^k\) by \eqref{tra:harmonic-overlaps:eq3}, division by that multiplicity proves \eqref{tra:harmonic-overlaps:eq13}.

\end{proof}

\subsection{Very small deficits by path coverage}

\begin{lemma}\label{tra:harmonic-overlaps:very-sparse}

\label{tra:harmonic-overlaps:very-sparse-statement}
There is one more preparatory bound. We can deal with \(1\le k=n-\lambda_1\le n^{0.51}\) directly using the binary layers in the sweep. For all sufficiently large \(n\) we will have
\begin{equation}
                         \|T_n\text{ in }\lambda\|_{\rm op}
                                             \le n^{-c k}                 \label{tra:harmonic-overlaps:eq27}
\end{equation}
for an absolute \(c>0\).

\end{lemma}\begin{proof}\label{tra:harmonic-overlaps:very-sparse-proof}

\label{tra:harmonic-overlaps:path-coverage-forest-proof}
We prove this on the harmonic part of \(\mathcal I_{n,k}\). Use distinct input and output tuples \(\xi,\zeta\). Given them a card has a unique required path through the sweep (it changes each bit to the output value when that bit is handled). For \(A\subseteq[k]\) write \(p_A\) for the probability the cards in \(A\) get the required outputs, and \(r_A=n^{|A|}p_A\). These probabilities are just obtained by imposing switch outcomes along paths. Two paths can pose conditions on a common switch only if they use the same switch in some layer. Call such paths adjacent, also when the needed outcomes may conflict. When the demands are compatible, the card paths being used cannot occupy the same wire at a time, and there is a factor of 2 in \(r_A\) for each shared switch, with no more than two cards sharing; a pair of cards in a compatible routing shares at most once (once they take opposite outputs on a switch, they stay different in that bit).

The centered kernel
\begin{equation}
                       \sum_{A\subseteq[k]}(-1)^{k-|A|} r_A              \label{tra:harmonic-overlaps:eq28}
\end{equation}
can be used instead of \(r_{[k]}\) on harmonic functions of \(\zeta\), since dropped-slot terms integrate to zero. This is for the slot transition forward on cards, acting on functions by output averaging (using an adjoint convention if necessary). The kernel \(r_{[k]}\) itself is the transition kernel against uniform on injections up to the factor from \(n^k\) versus \((n)_k\). Thus for \eqref{tra:harmonic-overlaps:eq27} it is enough to bound the Hilbert-Schmidt norm from \eqref{tra:harmonic-overlaps:eq28}, ignoring such a factor or allowing \(\exp(Ck)\). The centered kernel is zero unless every slot's path is adjacent to some other: isolated slots give cancellation since their normalized path factor is 1 and independent of the conditions on others.

It remains to bound the integrals of \(r_A^2\) on this coverage event for each \(A\), since this controls the squared norm up to \(\exp(Ck)\). We can do the integral against iid inputs and outputs restricted to distinct tuples, at cost \(\exp(Ck)\) by \eqref{tra:harmonic-overlaps:eq4}. One power of \(r_A\) gives a sampling using actual routed paths for the slots in \(A\): take a random full routing (choose all switches), and sample those input cards uniformly without replacement. Indeed under this sampling outputs for those cards have the probability \(p_A\) given their distinct inputs. Other slots can still just use iid input and output samples, ignoring distinctness for upper bounds. The factors from these measure comparisons can again be allowed as \(\exp(Ck)\). In this setting the remaining weight \(r_A\) on distinct tuples just counts factors 2 as above. Condition on the full routing. For probability bounds on constraints involving sampled paths, the actual routed paths can also be compared to sampling from them with replacement with a factor \(\exp(Ck)\), still evaluating any weight bound for distinct samples.

There are two useful observations here:
\paragraph{Conditional path adjacency.} Given a path of either type (from the actual full routing or from independent endpoints), a new independent choice of a path of either type, using with-replacement sampling conditional on the full routing, has probability at most \(2\log_2(n)/n\) to be adjacent to it. Per switch of the given path there are only two paths in the full routing using that switch, and for independently sampled endpoints the path goes through a uniform position at any fixed layer.
\paragraph{Shared-switch count.} For any \(b\) distinct members of a full routing the number of switches shared within that set is at most \((b/2)\log_2 b\). Indeed split inputs into halves according to the last bit handled in the sweep. Within halves before the last layer use induction, and for the last layer cross edges add at most the smaller half-count within the set. The bound follows since the binary entropy at any proportion is at least twice the smaller proportion.

For each component of the path adjacency graph of size \(b'\) the weight \(r_A\) therefore uses at most a factor \((b')^{b'/2}\), using observation 2 on actual slots. On the coverage event all components have size at least two. Use spanning forests of these components to bound the probabilities: with \(l\) total tree edges they can be counted with a bound \(2^{O(k)} k^l\) by orienting toward roots, and for any one forest the edge requirements cost probability at most \((2\log_2(n)/n)^l\), by observation 1 under with-replacement sampling. Here we may sum by forests describing the exact components from the event for which the weight bound is used, and just drop non-edge constraints for the probability upper bound.

These give a power saving \(n^{-c' k}\) even with the weights for components and all \(\exp(Ck)\) factors. Explicitly \(2k\log_2(n)/n\le n^{-0.47}\) for \(n\) large. In a component of size 2 we use one such factor and weight at most 2; in one of size \(b'\ge 3\) we use \(b'-1\) such factors, and weight bounded by \(n^{0.255 b'}\). This comparison applies uniformly in the forest count bound (multiply the maximum of each bound with \(l\) edges by the number of such forests). It proves the required Hilbert-Schmidt estimate and hence \eqref{tra:harmonic-overlaps:eq27}.

\end{proof}

\subsection{Interpolation through the split}

\begin{proof}[Proof of Theorem~\ref{tra:harmonic-overlaps:singular}]

\label{tra:harmonic-overlaps:deficit-ranges-and-dimensions}
We now combine the estimates. We only need to give the large-\(n\) induction step, with some absolute threshold for it chosen large enough throughout.

As before set \(L=\log n\). We organize the cases as follows apart from \eqref{tra:harmonic-overlaps:eq2}:
\begin{itemize}
\item Very small deficits as in \eqref{tra:harmonic-overlaps:eq27} can already use \eqref{tra:harmonic-overlaps:eq1} with any sufficiently small absolute \(a>0\), since, for example, \(D_\lambda\le n^k\) by occurring in the tuple representation.
\item For remaining shapes with \(k\le n\exp(-L^{1/4})\) use the sparse-deficit overlap estimate \eqref{tra:harmonic-overlaps:eq13}. Here \(k>n^{0.51}\), and
\end{itemize}
\begin{equation}
                          \log D_\lambda\ge c k L^{1/4}                  \label{tra:harmonic-overlaps:eq29}
\end{equation}
by \eqref{tra:harmonic-overlaps:eq3}.
\begin{itemize}
\item For the other shapes not killed by \eqref{tra:harmonic-overlaps:eq2}, we can use the regular trace overlap estimate \eqref{tra:harmonic-overlaps:eq5}; we have
\end{itemize}
\begin{equation}
                          \log D_\lambda\ge n\exp(-L^{1/3}).             \label{tra:harmonic-overlaps:eq30}
\end{equation}

To justify the last dimension assertion, we give the elementary checks. If \(\max(\lambda_1,\lambda'_1)< n/(4e)\), the hook bound gives exponential dimension since hook lengths are less than \(n/(2e)\). Otherwise take width \(w\ge n/(4e)\) using one of the shape or its transpose. If the number of remaining boxes is at least \(\lfloor w/8\rfloor\), we may just keep the full width and that many more boxes in a subshape; dimension is nondecreasing under adding boxes of a shape, and \eqref{tra:harmonic-overlaps:eq3} gives exponential dimension for the subshape. The case with fewer than that many remaining boxes in the transposed orientation would be killed by \eqref{tra:harmonic-overlaps:eq2}. With fewer in the original orientation we again get the claimed lower bound from \eqref{tra:harmonic-overlaps:eq3} when \(k> n\exp(-L^{1/4})\). This proves \eqref{tra:harmonic-overlaps:eq30}.

\label{tra:harmonic-overlaps:dyadic-analytic-family}
Suppose \eqref{tra:harmonic-overlaps:eq1} holds for the smaller sizes \(m,q\) with exponents at least \(a_*\), where without loss \(a_*\le 1/8\). In any row factor take its matrix singular decomposition in each shape, and similarly in columns. We need the bases meeting in multiplying across the two factors of \(T_n\). That is, removing outside unitaries, we bound singular values of the product of positive factors, using on one side left singular projections of its component operators and on the other right singular projections. This can all be done in the respective group algebras and then taken inside shape \(\lambda\), by ordinary group Fourier decomposition (the polar unitaries can be completed on nullspaces). Equivalently, in restriction of \(\lambda\) to the product subgroup, each component operator acts on its own irreducible by the corresponding matrix with identity on multiplicities.

In each component group and shape \(\mu\), group singular indices into dyadic blocks from \(r_0\) through at most \(2r_0-1\), \(r_0\) a power of 2. Use the projections onto these blocks in the meeting bases just described. Take \(R=D_\mu r_0\) as upper bound of regular rank in our estimates. Thus on the two sides we have orthogonal decompositions by the \(P_H,P_J\), respectively, with the regular rank product bounds \(R_H,R_J\); and the positive singular-value factors can use eigenvalue upper bounds
\begin{equation}
                                  R_H^{-a_*},\qquad R_J^{-a_*}             \label{tra:harmonic-overlaps:eq31}
\end{equation}
on the respective blocks. In fact for purposes of bounding any singular value of the product we can replace the factors by the indicated weights: individually the singular-value factors are then the weighted projection sums times commuting contractions, which can be taken to the outsides.

In shape \(\lambda\) consider the analytic matrix product of projection sums with weights, in the product order needed, i.e. of
\begin{equation}
                       \sum_{P_H} R_H^{-z\beta}P_H,\qquad
                       \sum_{P_J} R_J^{-z\beta}P_J,
             \qquad 0\le\operatorname{Re} z\le 1.                      \label{tra:harmonic-overlaps:eq32}
\end{equation}
Use either overlap estimate, with \(\beta\) chosen according to it:
\begin{itemize}
\item For shapes using \eqref{tra:harmonic-overlaps:eq5}, \(\beta=(1+\epsilon)/4\). On the line \(\operatorname{Re} z=1\), Schatten 4 norm of the product is bounded by
\end{itemize}
\begin{equation}
                         \left(\frac{\exp(E_n)}{D_\lambda}\right)^{1/4},
                    \quad E_n=n\exp(-L^{1/2})+ O(n^{4/5}).               \label{tra:harmonic-overlaps:eq33}
\end{equation}
Indeed expand in projection pairs and use the triangle inequality and \eqref{tra:harmonic-overlaps:eq5}, with regular trace paying for \(D_\lambda\) copies. The number of pairs can be absorbed with the additive \(O(n^{4/5})\) in the exponent. For a group of size-parameter \(M=m\) or \(q\), there are just \(\exp(O(\sqrt M))\) shapes and at most \(O(M\log M)\) dyad choices for each. We use here the usual partition count bound, also seen directly by using \(\exp(-1/\sqrt M)\) as argument in the partition generating function \(\prod_{i\ge 1}(1-x^i)^{-1}\), whose log at this argument is \(O(\sqrt M)\).
\begin{itemize}
\item For shapes using \eqref{tra:harmonic-overlaps:eq13} in the sparse range, take \(\beta=1/4\), and \eqref{tra:harmonic-overlaps:eq33} holds instead with exponent
\end{itemize}
\[
                                     E_n=O(k\log\log n).
\]
Here \eqref{tra:harmonic-overlaps:eq14} on both sides is necessary for nonzero terms, by restriction in the tuple calculation. Vectors of local deficits summing to at most \(k\) have at most \(\exp(O(k))\) choices since \(k>n^{0.51}\) exceeds both \(m,q\) for large \(n\). For positive deficit \(j\), shape choices cost at most \(2^j\), and dyad choices at most \(O(j\log n)\) by the dimension bound \(n^j\); for zero deficit there is just the trivial rank and shape. So the claimed exponent absorbs the number of terms and \eqref{tra:harmonic-overlaps:eq13}'s losses.

\label{tra:harmonic-overlaps:schatten-interpolation-and-uniform-exponent}
On \(\operatorname{Re} z=0\), \eqref{tra:harmonic-overlaps:eq32} gives operator norm at most 1 with no counting, by orthogonality within each of the two projection sums. We interpolate the whole product, which in particular allows counting costs in \eqref{tra:harmonic-overlaps:eq33} to be scaled by the interpolation parameter. By standard three-lines Schatten interpolation (between operator norm and Schatten 4), at parameter
\[
                                   \theta=a_*/\beta
\]
we obtain the bound \eqref{tra:harmonic-overlaps:eq33}, raised to power \(\theta\), in Schatten \(4/\theta\). The dual-matrix three-lines argument is the one given in the proof
of Lemma~\ref{lem:positive-power-comparison}; here its analytic
family is the full product \eqref{tra:harmonic-overlaps:eq32}.


This parameter gives just the weights \eqref{tra:harmonic-overlaps:eq31}, so ranking the resulting singular values at the root shows
\begin{equation}
                \log s_r(T_n\text{ in }\lambda)
                     \le -\frac{a_*}{4\beta}
                                         \big(\log(D_\lambda r)-E_n\big).    \label{tra:harmonic-overlaps:eq34}
\end{equation}
Use \eqref{tra:harmonic-overlaps:eq30} in the case using \eqref{tra:harmonic-overlaps:eq5}, or \eqref{tra:harmonic-overlaps:eq29} in the case using \eqref{tra:harmonic-overlaps:eq13}. With an absolute \(C\), say we conclude for all sufficiently large \(n\) the desired bound \eqref{tra:harmonic-overlaps:eq1} at the root with exponent at least
\begin{equation}
                                   a_*(1-C/L^{1/10}),                      \label{tra:harmonic-overlaps:eq35}
\end{equation}
provided we keep initial exponents small enough to allow the very small deficits as well. Crucially the choice of large-size threshold in these comparisons of exponents can be absolute and independent of how small \(a_*\) might be.

Finally these induction steps give \eqref{tra:harmonic-overlaps:eq1} with uniform positive exponent. Choose a large absolute base threshold, also large enough that the factors in \eqref{tra:harmonic-overlaps:eq35} are positive and, say, at least \(1/2\). At or below the threshold, Lemma~\ref{lem:finitegap} gives
\eqref{tra:harmonic-overlaps:eq1} with a common positive exponent,
chosen small enough for the very small deficits.
Lemma~\ref{lem:dyadic-exponents}, with $\gamma=1/10$, then bounds
the product of the losses in \eqref{tra:harmonic-overlaps:eq35}
away from zero. This proves \eqref{tra:harmonic-overlaps:eq1}.


\end{proof}

\subsection{Fourier summation and physical time}

\label{tra:harmonic-overlaps:fourier-completion}
This is the classical finite-group Fourier/Plancherel-to-total-variation
method of Diaconis and Shahshahani~\cite[Section 2, Lemma 2; Section 3,
proof of Lemma 14]{DiaconisShahshahani1981}. We use the noncentral matrix
form, not the random-transposition specialization; the harmonic lifting
and overlap estimates above are proved here.

Apply the indexed-bound conversion in
Section~\ref{sec:spectral-conversion} to
\eqref{tra:harmonic-overlaps:eq1}. After an integer $u$ of sweeps,
the chi-square divergence is at most
\begin{equation}
 \sum_{\lambda\ne(n),(1^n)}D_\lambda^{\,2-2au}.
 \label{tra:harmonic-overlaps:eq36}
\end{equation}
Sign and the other blocks in \eqref{tra:harmonic-overlaps:eq2}
vanish. If $2au\ge3$, Lemma~\ref{lem:degrees} makes the displayed
sum tend to zero. Thus a fixed number of sweeps gives vanishing
total-variation error, uniformly over the starting deck.
Lemma~\ref{lem:support} gives the lower bound $2d-O(1)$ in
physical shuffle units.
